\documentclass[../../main.tex]{subfiles}
\begin{document}

{\bf[解]}

与えられた不等式
\begin{align}
  a,b,c,d > 0 \label{eq:1}\\
  s(1-a)-tb>0 \label{eq:2}\\
  -sc+t(1-d)>0 \label{eq:3}
\end{align}
について考える．与えられた二次方程式の左辺を$f(x)$とおき：
\begin{align}
  f(x) = x^2-(a+d)x+(ad-bc)
\end{align}
$f(x)=0$について考える．

まずは\cref{eq:1,eq:2,eq:3}を満たす$t,s>0$が存在する$a,b,c,d$の条件について考える．\cref{eq:2,eq:3}より
\begin{align}
  \begin{cases}
    s(1-a) > tb \\
    t(1-d) > sc
  \end{cases}
\end{align}
だが，$t,b,s,c>0$よりこれらの右辺は正だから
\begin{align}
  \begin{cases}
    s(1-a)>0 \\
    t(1-d)>0
  \end{cases}
  \therefore
  \begin{cases}
    a<1 \\
    d<1
  \end{cases} \label{eq:4}
\end{align}
である．よって再度\cref{eq:2,eq:3}をこのもとで整理して
\begin{align}
  \dfrac{c}{1-d}s < t < \dfrac{1-a}{b}s
\end{align}
だから，このような$s,t>0$が存在する条件は
\begin{align}
  \dfrac{c}{1-d} < \dfrac{1-a}{b} \quad \therefore bc < (1-a)(1-d) \label{eq:5}
\end{align}
である．\cref{eq:4}と合わせて$a,b,c,d$の条件は
\begin{align}
  \begin{cases}
    0<a,d<1 \\
    0<b,c \\
    bc < (1-a)(1-d)
  \end{cases} \label{eq:6}
\end{align}
である．

\begin{figure}[htb]
  \centering
  \begin{tikzpicture}[scale=1.8, >=stealth, every node/.style={font=\small}]

    % --- 1. 座標定義 ---
    \coordinate (O) at (0, 0);

    % 直線の傾き: m1 = c/(1-d), m2 = (1-a)/b  (m1 < m2)
    % 例として m1 = 0.45, m2 = 1.15
    \def\mOne{0.45}
    \def\mTwo{1.15}
    \def\sMax{3.0}

    % --- 2. 満たすべき領域 (m1*s < t < m2*s, s > 0) のハッチング ---
    \fill[pattern=north east lines, pattern color=blue!65, opacity=0.7]
    (O) -- (\sMax, {\mOne*\sMax}) -- (\sMax, {\mTwo*\sMax}) -- cycle;

    % --- 3. 座標軸 ---
    \draw[->, thick, gray!80!black] (-0.3, 0) -- (3.5, 0) node[right] {$s$};
    \draw[->, thick, gray!80!black] (0, -0.3) -- (0, 3.8) node[above] {$t$};
    \node[below left=2pt] at (O) {$O$};

    % --- 4. 境界直線（境界は含まないので破線）---
    % (i) 下側境界: t = c/(1-d) s
    \draw[very thick, dashed, teal!85!black] (O) -- (3.2, {\mOne*3.2});
    \node[teal!90!black, font=\footnotesize, right] at (3.2, {\mOne*3.2})
    {$t = \frac{c}{1-d}s$};

    % (ii) 上側境界: t = (1-a)/b s
    \draw[very thick, dashed, blue!85!black] (O) -- (2.8, {\mTwo*2.8});
    \node[blue!90!black, font=\footnotesize, above left] at (2.7, {\mTwo*2.7})
    {$t = \frac{1-a}{b}s$};

    % 原点 (除外点: 白抜き丸)
    \draw[thick, fill=white] (O) circle (1.3pt);

  \end{tikzpicture}
  \caption{$(s,t)$ 平面において不等式を満たす正の点 $(s,t)$ の存在領域(斜線部境界を含まず)}
  \label{fig:st_plane_region}
\end{figure}

\cref{eq:6}の元で$f(x)=0$について考えると，判別式$D$として
\begin{align}
  \begin{cases}
    \text{端点}: f(\pm 1) = (1+ad-bc)\mp (a+d) >(a+d)\mp (a+d)>0 \\
    \text{軸}: -1 < \dfrac{a+d}{2} < 1 \\
    \text{判別式}: D = (a+d)^2-4(ad-bc) = (a-d)^2 + 4bc > 0
  \end{cases}
\end{align}
だから，$y=f(x)$のグラフの概形は下図のようになり，$f(x)=0$は$-1<x<1$に二つの異なる実数解を持つ．$\Box$

% TODO: y=f(x)の概形
\begin{figure}[htb]
  \centering
  \begin{tikzpicture}[scale=2.2, >=stealth, every node/.style={font=\small}]

    % --- 1. 定数・幾何パラメータ ---
    % 0 < a, d < 1, bc < (1-a)(1-d) の一例:
    % a = 0.5, d = 0.3 => a+d = 0.8 (軸 x_axis = 0.4)
    % ad = 0.15, bc = 0.10 => ad - bc = 0.05
    % f(x) = x^2 - 0.8x + 0.05
    % 頂点: (0.4, 0.05 - 0.16) = (0.4, -0.11)
    % f(1) = 1 - 0.8 + 0.05 = 0.25 > 0
    % f(-1) = 1 + 0.8 + 0.05 = 1.85 > 0
    % 解: x = 0.4 ± sqrt(0.16 - 0.05) = 0.4 ± sqrt(0.11) ≈ 0.4 ± 0.3317 (0.068, 0.732)
    \def\sumAD{0.8}
    \def\prodDet{0.05}
    \pgfmathsetmacro{\xAxis}{\sumAD / 2.0}                        % 0.4
    \pgfmathsetmacro{\yVertex}{\prodDet - \xAxis*\xAxis}          % -0.11
    \pgfmathsetmacro{\xSolOne}{\xAxis - sqrt(\xAxis*\xAxis - \prodDet)} % ≈ 0.068
    \pgfmathsetmacro{\xSolTwo}{\xAxis + sqrt(\xAxis*\xAxis - \prodDet)} % ≈ 0.732
    \pgfmathsetmacro{\yPosOne}{1.0 - \sumAD + \prodDet}           % 0.25
    \pgfmathsetmacro{\yNegOne}{1.0 + \sumAD + \prodDet}           % 1.85

    \coordinate (O)      at (0, 0);
    \coordinate (Vertex) at (\xAxis, \yVertex);
    \coordinate (P_neg1) at (-1.0, \yNegOne);
    \coordinate (P_pos1) at (1.0, \yPosOne);

    % --- 2. 区間 (-1, 1) のハイライト背景 ---
    \fill[blue!5, opacity=0.8] (-1.0, -0.4) rectangle (1.0, 2.1);

    % --- 3. 座標軸 ---
    \draw[->, thick, gray!80!black] (-1.4, 0) -- (1.5, 0) node[right] {$x$};
    \draw[->, thick, gray!80!black] (0, -0.4) -- (0, 2.2) node[above] {$y$};
    \node[below left=2pt] at (O) {$O$};

    % --- 4. 軸 x = (a+d)/2 の破線 ---
    \draw[dashed, teal!85!black, thick] (\xAxis, -0.35) -- (\xAxis, 1.9);
    \node[below=2pt, font=\scriptsize, teal!90!black] at (\xAxis, 0) {$\frac{a+d}{2}$};

    % --- 5. 放物線 y = f(x) のプロット ---
    \draw[line width=1.8pt, blue!85!black, domain=-1.1:1.35, samples=70]
    plot (\x, {\x*\x - \sumAD*\x + \prodDet});
    \node[blue!90!black, font=\footnotesize\bfseries, right] at (1.15, 1.25)
    {$y = f(x)$};

    % --- 6. 端点 x = ±1 における値 f(±1) ---
    % x = 1
    \draw[dashed, gray!60] (1.0, 0) -- (P_pos1) -- (0, \yPosOne);
    \fill[black] (P_pos1) circle (1.2pt);
    \node[below=2pt, font=\scriptsize] at (1.0, 0) {$1$};
    \node[left=2pt, font=\scriptsize, gray!90!black] at (0, \yPosOne) {$f(1) > 0$};

    % x = -1
    \draw[dashed, gray!60] (-1.0, 0) -- (P_neg1) -- (0, \yNegOne);
    \fill[black] (P_neg1) circle (1.2pt);
    \node[below=2pt, font=\scriptsize] at (-1.0, 0) {$-1$};
    \node[right=2pt, font=\scriptsize, gray!90!black] at (0, \yNegOne) {$f(-1) > 0$};

    % --- 7. 2つの実数解 (x 軸との交点) ---
    \fill[red!85!black] (\xSolOne, 0) circle (1.4pt);
    \fill[red!85!black] (\xSolTwo, 0) circle (1.4pt);

    % --- 8. 頂点 (y < 0) ---
    \fill[blue!90!black] (Vertex) circle (1.2pt);
    \node[below right=1pt, font=\scriptsize, blue!90!black] at (Vertex) {頂点 ($D > 0$)};

  \end{tikzpicture}
  \caption{放物線 $y = f(x)$ の概形と開区間 $(-1, 1)$ 内の 2 つの実数解}
  \label{fig:fx_two_real_roots}
\end{figure}

\newpage
\end{document}
